% Teil II: der durchgerechnete Weg.

% Gedankenschritt: Parität via 2R-n; Symmetrie nutzen, Gleichstand bei 0
% abspalten.

\subsection*{Gegeben und gesucht}
Start bei $0$, pro Sekunde $\pm1$ mit Wahrscheinlichkeit je $\tfrac12$.
Nach $n$ Sekunden: (a) mögliche Positionen für gerade $n$; (b) für
ungerade $n$; (c) Wahrscheinlichkeit für ``echt rechts von $0$``.

\subsection*{Ansatz}
Mit $R$ = Anzahl Rechtschritte, $L=n-R$ ist die Position $R-L=2R-n$.
Mögliche Positionen: $2R-n$ für $R=0,\dots,n$. Das ist eine reine
Linearkombination --- daher die Parität.

\subsection*{a) $n$ gerade}
Zahlwerte $2R-n$ sind gerade (gerade minus gerade). Alle geraden Zahlen
zwischen $-n$ und $n$:
\begin{equation}
  \{-n,\,-n+2,\dots,-2,\,0,\,2,\dots,n\} .
\end{equation}
Rückkehr zu $0$ ist möglich --- $R=\tfrac n2$ gibt $2\cdot\tfrac n2-n=0$.
\checkmark

\subsection*{b) $n$ ungerade}
Alle Positionen sind ungerade:
\begin{equation}
  \{-n,\,-n+2,\dots,-1,\,1,\dots,n\} .
\end{equation}
Hier ist $0$ \emph{nicht} erreichbar (kein ganzzahliges $R$ mit
$2R-n=0$).

\subsection*{c) Wahrscheinlichkeit ``echt rechts``}
$\mathrm P(R=r)=\binom{n}{r}2^{-n}$.
\begin{itemize}
  \item $n$ \emph{ungerade}: Symmetrie ($R\leftrightarrow L$) macht
        $P(\text{rechts})=P(\text{links})$, und es gibt keinen
        Gleichstand, die beiden decken also zusammen alle $n$-Schritt-
        Ergebnisse ab:
        \begin{equation}
          \mathrm P(\text{rechts})=\tfrac12 .
        \end{equation}
  \item $n$ \emph{gerade}: Der Gleichstand $R=n/2$ hat
        $\mathrm P(0)=\binom{n}{n/2}2^{-n}$. Über die Normierung
        $P(\text{rechts})+P(\text{links})+P(0)=1$ und die Symmetrie:
        \begin{equation}
          \mathrm P(\text{rechts})=\frac12\Bigl(1-\binom{n}{n/2}2^{-n}\Bigr).
        \end{equation}
\end{itemize}

\begin{ergebnis}
\begin{align*}
  \text{a)}&\ \{-n,-n+2,\dots,n\}\ \text{(nur gerade)}\\
  \text{b)}&\ \{-n,-n+2,\dots,n\}\ \text{(nur ungerade)}\\
  \text{c)}&\ P(\text{rechts})=
     \begin{cases}\frac12,& n\ \text{ungerade},\\
     \frac12\bigl(1-\binom{n}{n/2}2^{-n}\bigr),& n\ \text{gerade}.\end{cases}
\end{align*}
\end{ergebnis}

\begin{kontrolle}
\textbf{Numerische Stichprobe:} $n=10$: $\binom{10}{5}2^{-10}
=\frac{252}{1024}=0.246$, also $P(\text{rechts})=\frac12(1-0.246)
=0.377\approx0.377$; $n=11$: $P(\text{rechts})=0.5$. Für kleine $n$ direkt
über $\binom{n}{r}$ nachprüfbar.

\textbf{Grenzfall:} Für $n\to\infty$ geht $\binom{n}{n/2}2^{-n}\to0$
(Du Moivre-Laplace: der Gleichstands-Term verschwindet wie
$1/\sqrt{n}$), und beide Formeln laufen auf $\tfrac12$ zu --- der
Walk ``verliert`` keine systematische Drift. Das ist eine schöne
Konsistenz zwischen dem diskreten Zählergebnis und der grossen-$n$-
Intuition.
\end{kontrolle}

\begin{merksatz}
Die erreichbaren Positionen haben die Parität von $n$ (Position $=2R-n$);
die Rechts-Wahrscheinlichkeit ist bei ungeradem $n$ genau $\tfrac12$, bei
geradem $n$ um den Gleichstandsanteil $P(0)$ gedrückt.
\end{merksatz}
